% concatenated from jacobian-copy.tex
\documentclass{amsart}
%Additional packages
\usepackage{amssymb}% some symbols, as \nrightarrow, for instance
\usepackage{mathrsfs}% to produce caligraphic text: The foliation \F(f), for instance
%\usepackage{mathdots}%to make the reverse diagonal dots  
\usepackage{hyperref}
\usepackage{verbatim}
%\usepackage{colortbl}
%\usepackage{refcheck}
\usepackage{graphicx}
\usepackage{subfigure}
\usepackage{psfrag}
\usepackage{enumitem}

\newtheorem{theorem}{Theorem}[section]
\newtheorem{lemma}[theorem]{Lemma}
\newtheorem{corollary}[theorem]{Corollary}
\newtheorem{proposition}[theorem]{Proposition}
\newtheorem{conjecture}{Conjecture}
{\theoremstyle{definition}
\newtheorem{definition}[theorem]{Definition}
\newtheorem{example}[theorem]{Example}
\newtheorem{remark}{Remark}}

%Additional commands
\newcommand{\Z}{\mathbb{Z}}
\newcommand{\R}{\mathbb{R}}
\newcommand{\C}{\mathbb{C}}
\newcommand{\tatop}[2]{\genfrac{}{}{0pt}{10}{#1}{#2}}
\newcommand{\Res}{\textnormal{Res}}
\newcommand{\Root}{\textnormal{Root}}
%
%--------------------------------------------------------
\begin{document}
\title[On mappings with Jacobian 1]{On mappings with Jacobian one}


\author[Z. Jelonek]{Zbigniew Jelonek}

\address[Zbigniew Jelonek]{ Instytut Matematyczny, Polska Akademia Nauk, \'Sniadeckich 8, 00-656 Warszawa, Poland      
\newline  
              E-mail: {\tt najelone@cyf-kr.edu.pl}
              }

\subjclass[2010]{Primary: 14R15; Secondary: 14P99.}

\keywords{Jacobian conjecture, real Jacobian conjecture, local polynomial diffeomorphism}

\date{\today}

\thanks{The author is grateful to Professor Lev Birbrair for the helpful discussions}

\maketitle



\begin{abstract}
We show that the set $A(n, d)$ of polynomial automorphisms $F : \Bbb C^n \to \Bbb C^n$ of degree at most  $d$ and with $Jac(F ) = 1$ is Zariski closed.
In particular every irreducible component of the set $A(n,d)$ of polynomial mappings with Jacobian $1$ is either composed with polynomial automorphisms or 
(generically) with counterexamples to the Jacobian Conjecture. Moreover every such component has dimension at least $n^2-1.$ In particular
if the set $X(n,d)$ is irreducible, and $n\ge 3, d\ge 6$, then a generic element of this set is a counterexample to the Jacobian Conjecture.

\end{abstract}

\maketitle

\section{Introduction}
Nowdays there is a lot of interests in a counterexample to the Jacobian Conjecture. Here we try to explain 
 how possible counterexamples are located in the variety $X(n,d)$ of all polynomial mappings of degree bounded by $d$ and with Jacobian $1.$ It was suggested already by Dru\.zkowski and Rusek
 \cite{dr}, \cite{rus}, that the variety $X(n,d)$ of such polynomials mappings is not irreducible (in fact they have considered the special case of the variety of maps in the Dru\.zkowski
 form). Here we prove that every component of this set is either consists with polynomial automorphisms, either generically it is composed wiith  counterexamples 
to the Jacobian Conjecture. It does not mean that $X(n,d)$ is reducible, but if not, then a generic element of this set is a counterexample to the Jacobian Conjecture,
at least if $n\ge 3, d\ge 6.$

\section{Main Result}

\begin{lemma}
The set $A(n, d)$ of polynomial automorphisms $F : \Bbb C^n \to \Bbb C^n$ of degree at most  $d$ and with $Jac(F ) = 1$ is Zariski closed.
\end{lemma}

\begin{proof}
It is easy to see that the set $A(n, d)$ is a constructible subset of $\Omega(n,d)\cong\Bbb C^M$ (where  $\Omega(n,d)$ is a space of all  mappings $F:\Bbb C^n \to \Bbb C^n$ of degree at most $d$).  Indeed, consider the algebraic set 
$$S=\{ (F,G)\in \Omega(n,d)\times \Omega(n, d^{n-1}) ; F\circ G=identity\}.$$ Then the projection $\pi(S)$ of
$S$ on  $\Omega(n,d)$ is a constructible set by Chevaley theorem. But $A(n,d)=\pi(S) \cap \{ F : Jac(F)=1\},$ hence it is a constructible set.



Hence it is enough to prove that the set $A(n, d)$ is closed in the euclidian topology. Let $f_k \in A(n, d)$ and $f_k \to F.$ 
Hence $F = (F_1, ..., F_n)$ is a polynomial mapping with $Jac(F ) = 1.$ We can assume that $f_n (0) = 0.$ Indeed, $f_k (0) = b_k$ and $b_k \to b = F (0).$

 Now if we show that $f_k - b_k \to G \in A(n, d)$,  then also $f_k \to G + b.$ Now let $f_k(1, ..., 1) = (w_{k,1} , ..., w_{k,n}) = w_k \in \Bbb C^n.$
  Hence $w_k \to w = F (1, ..., 1).$  If $g_k = f^{ -1}_k,$ then $g_k(w_k) = (1, ..., 1).$  Note that the family of all $w_k$ is bounded.
  Now let $g_k = (Q_{k,1} , ..., Q_{k,n}).$ It is well known that the algebraic degree of $g_k$ is uniformly bounded by $d^{n-1}.$
  
   For a polynomial $Q = \sum a_\alpha x^\alpha\in  \Bbb C[x_1, ..., x_n]$ let $||Q|| =max |a_\alpha|.$ Let us denote $\tilde{Q} = Q/||Q||.$
    Since $Q_{k,j} (w_k) \to 1$ and $w_k$ is bounded, we have that there is an $\epsilon > 0$, such that $||Q_{k,j} || > \epsilon$ for every $k, j.$
    
    Now we have $Q_{k,1} (f_{k,1} , ..., f_{k,n} ) = x_1$. Consequently $\tilde{Q}_{k,1} (f_{k,1} , ..., f_{k,n} ) = a(k) x_1$, where $0 <a(k) < 1/\epsilon$. 
    If we take a subsequence of $\tilde{Q}_{k,1}$ we can assume that it converges to a polynomial $Q_1$ and one of its (fixed) coefficients has norm $1$. 
    Moreover we can assume that $a(k)$ also converges to $a$. Let us note that $a  \not= 0,$ because otherwise $Q_1(F_1, ..., F_n) = 0$ and $Q_1$ is a non-constant polynomial. 
    
    This contradicts the fact that $F_1, ..., F_n$ are algebraically independent, as components of a mapping with Jacobian equal to $1.$
     Hence finally $(1/a)Q_1(F_1, ..., F_n) = x_1,$ i.e., there is a polynomial $Q'_1$ such that $Q'_1(F_1, ..., F_n) = x_1.$
      In the same way there are polynomials $Q'_2, ..., Q'_n$ such that $Q'_j (F_1, ..., F_n) = x_j$ for $j = 2, ...n.$ This means that F is an automorphism.
\end{proof}

\begin{theorem} Let $X(n,d)$ denotes the space (which is an affine algebraic variety) of polynomial mappings $F$ with $Jac(F)=1$ and of degree bounded by $d.$ Let $Y$ be an irreducible component of this set. Then either $Y$ consists of polynomial automorphisms, or a generic element of $Y$ is a counterexample to Jacobian Conjecture. Moreover, $\dim Y \ge n^2-1.$
\end{theorem}

\begin{proof} Let $Y$ be an irreducible component. Then either $Y\subset A(n,d)$ or it contains a counterexample to the Jacobian Conjecture. In the latter case since $A(n,d)$ is closed it means that $A(n,d)\cap Y$ is a nowhere dense in $Y.$ Now denote by $\mathcal L$ the set of linear automorphisms with the Jacobian $1.$ Take $F\in Y$, which is not contaned in any different component of $X(n,d).$
Then $F\circ \mathcal{L}$ is an irreducible set which is contained in $X(n,d).$ Since it has a common point which belongs to $Y$ and does not belong to any other component it means that $F\circ \mathcal{L}\subset Y.$ But $\dim F\circ \mathcal{L}=n^2-1.$
\end{proof}

\begin{corollary}
If the set $X(n,d)$ is irreducible, and $n\ge 3, d\ge 6$, then a generic element of this set is a counterexample to the Jacobian Conjecture. 
\end{corollary}

\begin{proof}
Indeed, due to \cite{cjc} we know that $X(n,d)$ is not composed only with polynomial automorphisms.
\end{proof}



   

\begin{thebibliography}{99} 

\bibitem{dr} Drużkowski, L.
{\em An effective approach to Keller's Jacobian conjecture},
{\bf Math. Ann. 264},  303–313, (1983).

\bibitem{rus} Rusek, K.
{\em A geometric approach to Keller's Jacobian conjecture},
{\bf Math. Ann. 264}, 315–320, (1983).

\bibitem{cjc}  Alpöge, L.,  Hello there the Jacobian conjecture is false, X post, July 20, 2026.



    \end{thebibliography}
	

\end{document}

